% Conclusiones y recomendaciones de \nombreRody (rúbrica del Entregable 3: 4 puntos, «del alumno»).
% Escribirlas en primera persona, con tus palabras: 2 o 3 párrafos alcanzan. Reemplaza la línea
% \pendiente{...} de abajo por tu texto. Preguntas guía, no hace falta responderlas todas:
%   - ¿Qué aprendiste sobre concurrencia con este trabajo (worker pool, barrera, Spin, speedup)?
%   - ¿Qué resultado te sorprendió y por qué?
%   - ¿Qué limitación ves en el trabajo y qué recomendarías para mejorarlo o escalarlo?
%   - En lo que hiciste tú (tu PR), ¿qué harías distinto?
% Para compilar y ver el PDF: cd tp/informe/tp && ./compilar.sh  (sale en build/main.pdf).
\subsubsection*{\nombreRody}

Lo que más me dejó este trabajo es que en concurrencia el costo real no está donde uno lo busca.
Esperaba que el límite del speedup fuera una parte secuencial del algoritmo, y resultó ser el costo
de coordinar: repartir bloques y esperar la barrera en cada iteración pesa más a medida que agrego
workers. Eso explica por qué con ocho workers el speedup es \cifra{speedup-p8} y no ocho, y por
qué por debajo de unos \cifra{n-equilibrio}~viajes la versión concurrente no conviene. La
granularidad importó al revés de lo que suponía: trocear fino casi no cuesta, lo que castiga es
dejar workers sin trabajo.

La segunda lección fue sobre verificar. Un \texttt{errors: 0} de Spin no dice nada si el mismo
chequeo no es capaz de dar \texttt{errors: 1} cuando el algoritmo está mal. Construir un mutante
por propiedad me obligó a entender qué verifica exactamente cada búsqueda, y en ese proceso la
herramienta me dio tres resultados falsos sin avisar. Desde entonces desconfío de cualquier
verificación que nunca vi fallar, sea un test, un modelo o una medición.

Como recomendación, mantendría el modelo Promela junto al código y en la integración continua, y
modelaría el cierre del pool antes de reutilizarlo en un programa de larga vida. Para escalar,
extendería la medición a varios meses de datos, donde coordinar pesa menos frente al trabajo, y
para pasar a varias máquinas repartiría los bloques por red manteniendo la reducción en orden fijo,
que es lo que hace reproducible el resultado.
